% Generated draft: requires mathematical and source review.
\documentclass[11pt]{article}
\usepackage[margin=25mm]{geometry}
\usepackage{fontspec}
\setmainfont{Latin Modern Roman}
\usepackage{amsmath,amssymb,mathtools,unicode-math}
\setmathfont{Latin Modern Math}
\usepackage{enumitem,needspace,xurl,hyperref}
\setlength{\parindent}{0pt}
\setlength{\parskip}{5pt}
\setlength{\emergencystretch}{3em}
\newcommand{\N}{\mathbb{N}}
\newcommand{\Z}{\mathbb{Z}}
\newcommand{\Q}{\mathbb{Q}}
\newcommand{\R}{\mathbb{R}}
\newcommand{\C}{\mathbb{C}}
\newcommand{\F}{\mathbb{F}}
\newcommand{\A}{\mathbb{A}}
\newcommand{\PP}{\mathbb{P}}
\newcommand{\E}{\mathbb{E}}
\newcommand{\eps}{\varepsilon}
\newcommand{\dd}{\,\mathrm{d}}
\begin{document}
\section*{UnsolvedMath ID: 30001258}
\textbf{Source problem code:} catalogue entry
\subsection{Definitions and mathematical statement}
Let \(\Omega\subset\R^N\) be a bounded domain and let \(1<q<p\), so that \(p>1\). The dimension \(N\), further hypotheses on \(\Omega\), and the admissible values and dependence on \(q\) of the parameter \(\lambda\) are not specified in the supplied record. The \(p\)-Laplacian is
\[
\Delta_p u=\operatorname{div}(|\nabla u|^{p-2}\nabla u).
\]
The stated Lane--Emden boundary value problem is
\[
\begin{cases}
-\Delta_pu=\lambda |u|^{q-2}u & \text{in }\Omega,\\
u=0 & \text{on }\partial\Omega.
\end{cases}
\]
A standard weak interpretation of the boundary condition is \(u\in W^{1,p}_0(\Omega)\), with
\[
\int_\Omega |\nabla u|^{p-2}\nabla u\cdot\nabla\varphi\,\dd x
=\lambda\int_\Omega |u|^{q-2}u\varphi\,\dd x
\qquad\text{for every }\varphi\in W^{1,p}_0(\Omega).
\]
Here \(|u|^{q-2}u\) is interpreted as zero at \(u=0\). This weak formulation is an explicit interpretive framework; the supplied record does not specify the intended solution regularity.

The associated variational functional is
\[
J_{\lambda,q}(u)=\frac1p\int_\Omega |\nabla u|^p\,\dd x
-\frac{\lambda}{q}\int_\Omega |u|^q\,\dd x.
\]
Under the usual least-energy convention, a ground state is a nonzero weak solution minimizing \(J_{\lambda,q}\) among nonzero weak solutions, and a least-energy nodal solution minimizes it among weak solutions whose positive and negative parts \(u^+=\max\{u,0\}\) and \(u^-=\max\{-u,0\}\) are both nonzero. These conventions make the terminology explicit but are not definitions confirmed by the supplied source; existence or attainment of either minimum is not asserted.

\textbf{Question as supplied.} Do the asymptotic results for ground states and least-energy nodal solutions of this problem as \(q\to p\) remain valid in the range \(1<q<p\), hence as \(q\uparrow p\)?

The supplied record does not state the asymptotic results to be extended, including their hypotheses, normalization, convergence topology, or limiting objects. Consequently, it does not determine a precise asymptotic assertion to prove or disprove. Its dated literature assessment also refers to an unestablished full \(q>p\) asymptotic picture, whereas the extracted question asks about extension to \(q<p\); this relationship remains unresolved.
\subsection{Short English statement}
The question asks whether unspecified asymptotic results for ground states and least-energy sign-changing solutions of the Dirichlet Lane--Emden problem extend to \(q\uparrow p\) with \(1<q<p\). The exact conclusions, parameter assumptions, and solution conventions require clarification before a precise theorem can be formulated.
\subsection{Sources}
\begin{itemize}
\item Supplied local registry record: permanent UnsolvedMath ID \(30001258\), title \emph{Sublinear Lane--Emden Ground-State Asymptotics}, exact source problem code \texttt{catalogue entry}, in the Oberwolfach Reports Open Problems collection. The proposed year is 2009; no individual proposer is supplied.
\item \emph{Topological and Variational Methods for Partial Differential Equations} (2009), Oberwolfach Reports, pp.\ 1427--1486, DOI 10.4171/owr/2009/26. Supplied report citation and literature-source link: \url{https://doi.org/10.4171/owr/2009/26}. The background describes this report as documenting workshop discussions and open problems; no precise problem page or speaker attribution is supplied.
\item Supplied literature-source label: ``inverse iteration for sublinear p-Laplacian problems,'' \url{https://arxiv.org/abs/1011.3172}. No authors or more detailed bibliographic attribution are supplied.
\item Literature assessment checked at \texttt{2026-08-21}: the supplied curation labels the problem partially solved and states that, for \(q<p\), existence and uniqueness results and \(p\)-Laplacian eigenpair limits are known, while the full \(q>p\) asymptotic picture requested in the extracted question has not been established. This is supplied dated curation, not an independently verified literature assessment; its hypotheses, scope, and relationship to the extracted question require review.
\item The registry labels the statement \texttt{corrected\_verified}; its accompanying verification note concerns balanced mathematical delimiters only. It does not establish mathematical correctness or verify the literature claims.
\end{itemize}
\end{document}
